\renewcommand{\bibname}{Bibliografía y recursos verificables}
\begin{thebibliography}{99}
\addcontentsline{toc}{chapter}{Bibliografía y recursos verificables}
\item[] Todas las páginas web se consultaron para el corte del 1 de octubre de 2026. Una URL de rama principal puede cambiar; únicamente se presenta como snapshot el commit indicado explícitamente. Las referencias no son una clasificación de calidad ni una validación independiente de las afirmaciones de rendimiento de sus autores.
\bibitem{tinygrad-pin} tinygrad. Snapshot de código, commit \texttt{c3aec477b99d9\allowbreak{}bb87c54d91897\allowbreak{}cf60acd3f17441}, 30-09-2026. \url{https://github.com/tinygrad/tinygrad/tree/c3aec477b99d9bb87c54d91897cf60acd3f17441}.
\bibitem{tinygrad-dev} tinygrad. \emph{Developer documentation}. \url{https://docs.tinygrad.org/developer/developer/}.
\bibitem{tinygrad-site} Tiny Corp. Sitio y explicación de operaciones. \url{https://tinygrad.org/}.
\bibitem{tinygrad-codegen} tinygrad. Pipeline de generación del snapshot. \url{https://github.com/tinygrad/tinygrad/blob/c3aec477b99d9bb87c54d91897cf60acd3f17441/tinygrad/codegen/__init__.py}.
\bibitem{tinygrad-org} Tiny Corp. Organización pública y repositorios. \url{https://github.com/tinygrad/}.
\bibitem{geohot} George Hotz. \emph{Five years of tinygrad}, 29-12-2025. \url{https://geohot.github.io/blog/jekyll/update/2025/12/29/five-years-of-tinygrad.html}.
\bibitem{cuda-wsl} NVIDIA. \emph{CUDA on WSL User Guide}. \url{https://docs.nvidia.com/cuda/wsl-user-guide/index.html}.
\bibitem{ieee-nvidia} NVIDIA. \emph{Floating Point and IEEE 754 Compliance for NVIDIA GPUs}. \url{https://docs.nvidia.com/cuda/floating-point/index.html}.
\bibitem{cuda} NVIDIA. \emph{CUDA Programming Guide}. \url{https://docs.nvidia.com/cuda/cuda-programming-guide/}.
\bibitem{ptx} NVIDIA. \emph{Parallel Thread Execution ISA}. \url{https://docs.nvidia.com/cuda/parallel-thread-execution/index.html}.
\bibitem{cutlass} NVIDIA. CUTLASS, README y código; sección 4.8.0 de septiembre de 2026. \url{https://github.com/NVIDIA/cutlass}.
\bibitem{cutile} NVIDIA. cuTile Python. \url{https://github.com/NVIDIA/cutile-python}.
\bibitem{nvidia-driver} NVIDIA. Open GPU Kernel Modules. \url{https://github.com/NVIDIA/open-gpu-kernel-modules}.
\bibitem{hipblaslt} AMD. hipBLASLt, aviso de migración al monorepo. \url{https://github.com/ROCm/hipBLASLt}.
\bibitem{ck} AMD. Composable Kernel, aviso de migración. \url{https://github.com/ROCm/composable_kernel}.
\bibitem{rocm-libraries} AMD. ROCm Libraries. \url{https://github.com/ROCm/rocm-libraries}.
\bibitem{aiter} AMD. AITER. \url{https://github.com/ROCm/aiter}.
\bibitem{rocm-llvm} AMD. ROCm LLVM Project. \url{https://github.com/ROCm/llvm-project}.
\bibitem{hip} AMD. HIP documentation. \url{https://rocm.docs.amd.com/projects/HIP/en/latest/}.
\bibitem{linux-amdgpu} Linux kernel documentation. AMDGPU driver. \url{https://docs.kernel.org/gpu/amdgpu/index.html}.
\bibitem{linux-mm} Linux kernel documentation. DRM memory management. \url{https://docs.kernel.org/gpu/drm-mm.html}.
\bibitem{linux-source} Linux. Fuente del kernel. \url{https://github.com/torvalds/linux}.
\bibitem{deepgemm} DeepSeek. DeepGEMM. \url{https://github.com/deepseek-ai/DeepGEMM}.
\bibitem{deepgemm-ascend} DeepSeek. DeepGEMM-Ascend. \url{https://github.com/deepseek-ai/DeepGEMM-Ascend}.
\bibitem{deepep} DeepSeek. DeepEP. \url{https://github.com/deepseek-ai/DeepEP}.
\bibitem{deepep-ascend} DeepSeek. DeepEP-Ascend: README y restricciones de la entrega preliminar. \url{https://github.com/deepseek-ai/DeepEP-Ascend}.
\bibitem{deepjit} DeepSeek. DeepJIT. \url{https://github.com/deepseek-ai/DeepJIT}.
\bibitem{ascend-samples} Huawei Ascend. Samples. \url{https://github.com/Ascend/samples}.
\bibitem{tilelang} Tile-AI. TileLang. \url{https://github.com/tile-ai/tilelang}.
\bibitem{tilelang-doc} TileLang. Documentación y API. \url{https://www.tilelang.com/autoapi/}.
\bibitem{tilelang-paper} Wang et al. \emph{TileLang: A Composable Tiled Programming Model for AI Systems}. 2025. \url{https://arxiv.org/abs/2504.17577}.
\bibitem{triton-tutorial} Triton. Tutorial de multiplicación de matrices. \url{https://triton-lang.org/main/getting-started/tutorials/03-matrix-multiplication.html}.
\bibitem{triton} Triton. Compilador y lenguaje. \url{https://github.com/triton-lang/triton}.
\bibitem{mlir-toy} LLVM. MLIR Toy Tutorial. \url{https://mlir.llvm.org/docs/Tutorials/Toy/}.
\bibitem{mlir-paper} Lattner et al. \emph{MLIR: A Compiler Infrastructure for the End of Moore's Law}. 2020. \url{https://arxiv.org/abs/2002.11054}.
\bibitem{llvm-langref} LLVM. Language Reference Manual. \url{https://llvm.org/docs/LangRef.html}.
\bibitem{tvm} Chen et al. \emph{TVM: An Automated End-to-End Optimizing Compiler for Deep Learning}. 2018. \url{https://arxiv.org/abs/1802.04799}.
\bibitem{egg} Willsey et al. \emph{egg: Fast and Extensible Equality Saturation}. Trabajo original POPL 2021; presentación CACM 2026. \url{https://doi.org/10.1145/3815481}.
\bibitem{roofline} Williams, Waterman y Patterson. \emph{Roofline: An Insightful Visual Performance Model for Multicore Architectures}. 2009. \url{https://doi.org/10.1145/1498765.1498785}.
\bibitem{attention} Vaswani et al. \emph{Attention Is All You Need}. 2017. \url{https://arxiv.org/abs/1706.03762}.
\bibitem{flash1} Dao et al. \emph{FlashAttention: Fast and Memory-Efficient Exact Attention with IO-Awareness}. 2022. \url{https://arxiv.org/abs/2205.14135}.
\bibitem{flash2} Dao. \emph{FlashAttention-2: Faster Attention with Better Parallelism and Work Partitioning}. 2023. \url{https://arxiv.org/abs/2307.08691}.
\bibitem{flash3} Shah et al. \emph{FlashAttention-3: Fast and Accurate Attention with Asynchrony and Low-precision}. 2024. \url{https://arxiv.org/abs/2407.08608}.
\bibitem{flash-repo} Dao-AILab. FlashAttention: implementaciones y requisitos por backend, incluida FA4 con CuTe DSL. \url{https://github.com/Dao-AILab/flash-attention}.
\bibitem{thunderkittens} HazyResearch. ThunderKittens. \url{https://github.com/HazyResearch/ThunderKittens}.
\bibitem{llmc} Karpathy et al. llm.c. \url{https://github.com/karpathy/llm.c}.
\bibitem{nanogpt} Karpathy. nanoGPT. \url{https://github.com/karpathy/nanoGPT}.
\bibitem{kernelbench} ScalingIntelligence. KernelBench. \url{https://github.com/ScalingIntelligence/KernelBench}.
\bibitem{deepseek-v3} DeepSeek-AI. \emph{DeepSeek-V3 Technical Report}. 2024. \url{https://arxiv.org/abs/2412.19437}.
\bibitem{deepseek-v4} DeepSeek-AI. DeepSeek-V4-Pro, model card y reporte enlazado. \url{https://huggingface.co/deepseek-ai/DeepSeek-V4-Pro}. Reporte: \url{https://arxiv.org/abs/2606.19348}.
\bibitem{taic} Costa et al. \emph{AI as a Compiler: Compiling Triton kernels without the Triton compiler}. 29-09-2026. \url{https://arxiv.org/abs/2609.36800}.
\bibitem{tilesight} Mo et al. \emph{TileSight: A First-Principles Tile-Centric Analytical GPU Performance Model from Cores to Clusters}. 2026. \url{https://arxiv.org/abs/2607.22432}.
\bibitem{correct-slow} \emph{Correct but Slow: An Empirical Study of the GPU Kernel Evaluation Gap in Modern Domain-Specific Languages}. 2026. \url{https://arxiv.org/abs/2607.04454}.
\bibitem{tile-bugs} \emph{Characterizing Real-World Bugs in Tile Programs for Automated Bug Detection}. 2026. \url{https://arxiv.org/abs/2605.19652}.
\bibitem{cuda-gpus} NVIDIA. CUDA GPU Compute Capability. \url{https://developer.nvidia.com/cuda/gpus}.

\bibitem{adam} D. P. Kingma y J. Ba. \emph{Adam: A Method for Stochastic Optimization}. 2014, ICLR 2015. \url{https://arxiv.org/abs/1412.6980}.
\bibitem{adamw} I. Loshchilov y F. Hutter. \emph{Decoupled Weight Decay Regularization}. ICLR 2019. \url{https://arxiv.org/abs/1711.05101}.
\bibitem{gqa} J. Ainslie y colaboradores. \emph{GQA: Training Generalized Multi-Query Transformer Models from Multi-Head Checkpoints}. 2023. \url{https://arxiv.org/abs/2305.13245}.
\bibitem{rope} J. Su y colaboradores. \emph{RoFormer: Enhanced Transformer with Rotary Position Embedding}. 2021. \url{https://arxiv.org/abs/2104.09864}.
\bibitem{rmsnorm} B. Zhang y R. Sennrich. \emph{Root Mean Square Layer Normalization}. 2019. \url{https://arxiv.org/abs/1910.07467}.
\bibitem{swiglu} N. Shazeer. \emph{GLU Variants Improve Transformer}. 2020. \url{https://arxiv.org/abs/2002.05202}.


\bibitem{zero} S. Rajbhandari y colaboradores. \emph{ZeRO: Memory Optimizations Toward Training Trillion Parameter Models}. 2019/2020. \url{https://arxiv.org/abs/1910.02054}.
\bibitem{megatron} M. Shoeybi y colaboradores. \emph{Megatron-LM: Training Multi-Billion Parameter Language Models Using Model Parallelism}. 2019/2020. \url{https://arxiv.org/abs/1909.08053}.
\bibitem{paged} W. Kwon y colaboradores. \emph{Efficient Memory Management for Large Language Model Serving with PagedAttention}. SOSP 2023. \url{https://arxiv.org/abs/2309.06180}.
\bibitem{chinchilla} J. Hoffmann y colaboradores. \emph{Training Compute-Optimal Large Language Models}. 2022. \url{https://arxiv.org/abs/2203.15556}.
\bibitem{nccl} NVIDIA. NCCL. \url{https://github.com/NVIDIA/nccl}.
\bibitem{rccl} AMD. RCCL: aviso de migración a ROCm/rocm-systems. \url{https://github.com/ROCm/rccl}.

\bibitem{tinyos} Tiny Corp. tinyos: constructor de imágenes de sistema para tinybox. \url{https://github.com/tinygrad/tinyos}.
\bibitem{tiny-mec} Tiny Corp. custom\_mec\_public: firmware MEC, emulador y pruebas para gfx1100. \url{https://github.com/tinygrad/custom_mec_public}.
\end{thebibliography}
